\documentclass[12pt]{article}
\usepackage[utf8]{inputenc}
\usepackage[spanish]{babel}
\usepackage[a4paper, margin=2cm]{geometry}
\usepackage{tikz}
\usetikzlibrary{babel}
\usepackage[most]{tcolorbox}
\usepackage{amsmath}

% Usar fuente sin serifa para un aspecto más moderno (tipo infografía)
\renewcommand{\familydefault}{\sfdefault}

% Definición de colores con temática médica/científica
\definecolor{medblue}{RGB}{43, 108, 176}
\definecolor{medteal}{RGB}{49, 151, 149}
\definecolor{medlight}{RGB}{235, 248, 255}
\definecolor{meddark}{RGB}{26, 32, 44}
\definecolor{bloodred}{RGB}{220, 38, 38}   % Rojo intenso para el glóbulo
\definecolor{bloodlight}{RGB}{254, 226, 226} % Rojo muy claro para el relleno

\begin{document}
\pagestyle{empty}

\begin{tcolorbox}[
    enhanced, 
    colback=medlight!30, 
    colframe=medblue, 
    title={\Large \textbf{Potencias de Diez: Notación Científica}}, 
    halign title=center, 
    fonttitle=\bfseries, 
    arc=4mm, 
    boxrule=1.5pt,
    shadow={2mm}{-2mm}{0mm}{black!30!white}
]

    \vspace{0.2cm}
    \begin{tcolorbox}[colback=white, colframe=medteal, arc=2mm, boxrule=1.2pt, title=\textbf{Caso Clínico / Problema}]
        Un glóbulo rojo (eritrocito) humano tiene un diámetro aproximado de \textbf{0.0000075} metros. Exprese esta medida utilizando notación científica (potencias de diez) y el prefijo adecuado del Sistema Internacional.
    \end{tcolorbox}

    \vspace{0.4cm}
    
    \begin{center}
    \begin{tikzpicture}[scale=1]
        % Eritrocito (vista superior simulando disco bicóncavo)
        \filldraw[bloodlight, draw=bloodred, line width=2.5pt] (0,0) circle (2cm);
        \filldraw[bloodred!20, draw=bloodred!50, line width=1pt] (0,0) circle (1.2cm);
        \filldraw[bloodred!40, draw=none] (0,0) circle (0.6cm); % Sombra central más profunda
        
        % Etiqueta del Eritrocito
        \node[font=\bfseries, text=bloodred!80!black] at (0, 0) {Eritrocito};
        
        % Acotación del diámetro
        \draw[<->, thick, meddark] (-2, -2.6) -- (2, -2.6) node[midway, below, font=\bfseries] {$d = 0.0000075 \text{ m}$};
        \draw[dashed, meddark!60, thick] (-2, -2.4) -- (-2, -0.5);
        \draw[dashed, meddark!60, thick] (2, -2.4) -- (2, -0.5);
        
        % Representación visual del salto de la coma decimal
        \node[font=\Large, text=meddark] at (0, -4.8) {
            $0 \textcolor{bloodred}{\mathbf{,}} \underbrace{0 \quad 0 \quad 0 \quad 0 \quad 0 \quad 7}_{\text{\small \textbf{6 lugares a la derecha}}} \textcolor{bloodred}{\mathbf{,}} 5 \text{ m} \quad \xrightarrow{\text{Notación}} \quad \mathbf{7.5 \times 10^{-6} \text{ m}}$
        };
        
    \end{tikzpicture}
    \end{center}

    \vspace{0.4cm}
    \begin{tcolorbox}[colback=medteal!10, colframe=medteal, arc=2mm, boxrule=1.2pt, title=\textbf{Desarrollo Matemático y Solución}]
        Para convertir la medida a notación científica, la coma decimal se desplaza \textbf{6 posiciones} hacia la derecha, obteniendo un coeficiente entre 1 y 10, multiplicado por una base 10 con exponente negativo:
        \[ d = \mathbf{7.5 \times 10^{-6} \text{ m}} \]
        En el Sistema Internacional (SI), la potencia $10^{-6}$ corresponde al prefijo \textbf{micro} ($\mu$). Sustituimos la potencia por el símbolo del prefijo:
        \[ d = \mathbf{7.5 \, \mu\text{m}} \]
        \textbf{Resultado:} El diámetro del eritrocito es \textbf{$7.5 \times 10^{-6}$ m}, lo que equivale a \textbf{7.5 micrómetros ($\mu$m)}.
    \end{tcolorbox}
    
\end{tcolorbox}

\end{document}